\subsection{Five click reactions as typed reductions}
\label{sec:click-chem}

Section~\ref{sec:mlc-beta} introduced the rule combinator
$\mathsf{R}_{k}$ as a placeholder for \emph{any} click reaction.
This subsection unpacks the placeholder: it lists the five click
reactions that the shipped implementation actually exposes,
writes each as a typed reduction, and reports an ablation on the
relative contribution of each reaction to search diversity.

\paragraph{Honest framing.}
All five rule definitions, SMARTS templates, and catalyst labels
are \textsc{measured} -- imported verbatim from
\texttt{molmetal/molmetal\_lam/lam\_chem/rules.py} (the
public-facing registry) which re-exports the
\texttt{ReactionRule} singletons from
\texttt{molmetal/molmetal\_lam/reactions/beta\_reductions.py}.
The ablation result is \textsc{measured} on
\texttt{wf\_lambda1\_build}. The five rules cover the
\emph{click-chemistry} subset of the reaction space; full
coverage of general synthetic chemistry is \textsc{projected}
future work (Section~\ref{sec:future}).

\subsubsection{The five reactions}
\label{sec:click-chem-five}

Table~\ref{tab:click-rules} lists the five rule combinators
$\mathsf{R}_{k}$ shipped in \texttt{rules.py}, with the
\texttt{ReactionRule} subclass that implements each and the
source-line anchor in
\texttt{molmetal/molmetal\_lam/reactions/beta\_reductions.py}.

\begin{table}[h]
  \centering
  \caption{The five click-reaction rule combinators
  $\mathsf{R}_{k}$. All five are imported verbatim from
  \texttt{molmetal/molmetal\_lam/reactions/beta\_reductions.py};
  the SMARTS templates are reproduced without modification.}
  \label{tab:click-rules}
  \begin{tabular}{lll}
    \hline
    Combinator $\mathsf{R}_{k}$ & Reaction (substrates $\to$ product) & Source line \\
    \hline
    $\mathsf{R}_{\text{CuAAC}}$         & azide + terminal alkyne $\to$ 1,4-disubstituted 1,2,3-triazole          & \texttt{beta\_reductions.py:563} \\
    $\mathsf{R}_{\text{SPAAC}}$         & azide + cyclooctyne $\to$ triazole                                       & \texttt{beta\_reductions.py:627} \\
    $\mathsf{R}_{\text{ThiolEne}}$      & thiol + alkene $\to$ thioether                                           & \texttt{beta\_reductions.py:817} \\
    $\mathsf{R}_{\text{Suzuki}}$        & aryl-boronic acid + aryl-halide $\to$ biaryl                             & \texttt{beta\_reductions.py:965} \\
    $\mathsf{R}_{\text{AmideCoupling}}$ & carboxylic-acid + amine $\to$ amide                                      & \texttt{beta\_reductions.py:1033} \\
    \hline
  \end{tabular}
\end{table}

\subsubsection{Typed-reaction signatures}
\label{sec:click-chem-types}

Each $\mathsf{R}_{k}$ carries a $\Pi$-typed signature: the rule
takes two typed substrate variables and returns a typed product.
Concretely, the rule combinators introduced in
\S\ref{sec:mlc-types} instantiate as
\begin{equation}
  \label{eq:cuaac-typed}
  \mathsf{R}_{\text{CuAAC}}
    \;:\;
    \Pi(\,a : \tau_{\text{azide}},\;
        b : \tau_{\text{alkyne}}\,).\;
        \tau_{\text{1,4-triazole}}
\end{equation}
\begin{equation}
  \label{eq:spaac-typed}
  \mathsf{R}_{\text{SPAAC}}
    \;:\;
    \Pi(\,a : \tau_{\text{azide}},\;
        b : \tau_{\text{cyclooctyne}}\,).\;
        \tau_{\text{triazole}}
\end{equation}
\begin{equation}
  \label{eq:thiolene-typed}
  \mathsf{R}_{\text{ThiolEne}}
    \;:\;
    \Pi(\,a : \tau_{\text{thiol}},\;
        b : \tau_{\text{alkene}}\,).\;
        \tau_{\text{thioether}}
\end{equation}
\begin{equation}
  \label{eq:suzuki-typed}
  \mathsf{R}_{\text{Suzuki}}
    \;:\;
    \Pi(\,a : \tau_{\text{aryl-B(OH)$_2$}},\;
        b : \tau_{\text{aryl-X}}\,).\;
        \tau_{\text{biaryl}}
\end{equation}
\begin{equation}
  \label{eq:amide-typed}
  \mathsf{R}_{\text{AmideCoupling}}
    \;:\;
    \Pi(\,a : \tau_{\text{carboxylic-acid}},\;
        b : \tau_{\text{amine}}\,).\;
        \tau_{\text{amide}}
\end{equation}
and each rule fires as a typed $\beta$-redex
\begin{equation}
  \label{eq:click-redex}
  \mathsf{R}_{k}\;a\;b
    \;\longrightarrow_{\beta}\;
    t
\end{equation}
where $t$ is the typed product combinator. CuAAC therefore
expands to the explicit term
\begin{equation}
  \label{eq:cuaac-reduction}
  (\mathsf{R}_{\text{CuAAC}}
   \;:\;
   \Pi(\tau_{\text{azide}},\, \tau_{\text{alkyne}})
     .\, \tau_{\text{1,4-triazole}})
   \;\cdot\;
   (\mathsf{C}_{a}^{\Sigma_{\text{azide}}})
   \;\cdot\;
   (\mathsf{C}_{b}^{\Sigma_{\text{alkyne}}})
   \;\longrightarrow_{\beta}\;
   \mathsf{C}_{t}^{\Sigma_{\text{1,4-triazole}}}.
\end{equation}
The four remaining reactions are notated identically; only the
$\Pi$-signature and the product combinator change.

\subsubsection{Verbatim SMARTS templates}
\label{sec:click-chem-smarts}

Each rule carries a SMARTS template (taken verbatim from the
source -- \textsc{measured}, not paraphrased) that the
$\beta$ reduction engine applies at the leftmost-outermost
substrate pair during $\beta$-reduction. The five templates are
\begin{align}
  \mathsf{R}_{\text{CuAAC}}\;:\;&
    \texttt{[N:1]=[N:2]=[N:3].[C:4]\#[CH:5]>>[C:4]1=[C:5][N:3]=[N:2][N:1]1}
  \tag{CuAAC}
  \label{eq:smarts-cuaac}\\
  \mathsf{R}_{\text{SPAAC}}\;:\;&
    \texttt{[N:1]=[N:2]=[N:3].[C:4]\#[C:5]>>[C:4]1=[C:5][N:3]=[N:2][N:1]1}
  \tag{SPAAC}
  \label{eq:smarts-spaac}\\
  \mathsf{R}_{\text{ThiolEne}}\;:\;&
    \texttt{[\#16:1][SH].[\#6:2]=[\#6:3]>>[\#16:1][\#6:2][\#6:3]}
  \tag{ThiolEne}
  \label{eq:smarts-thiolene}\\
  \mathsf{R}_{\text{Suzuki}}\;:\;&
    \texttt{[\#6:1][B]([O])[O].[\#6:3][F,Cl,Br,I]>>[\#6:1][\#6:3]}
  \tag{Suzuki}
  \label{eq:smarts-suzuki}\\
  \mathsf{R}_{\text{AmideCoupling}}\;:\;&
    \texttt{[C:1](=[O:2])[OH].[NH2:4]>>[C:1](=[O:2])[NH:4]}
  \tag{AmideCoupling}
  \label{eq:smarts-amide}
\end{align}
All five templates parse under RDKit's SMARTS parser and the
product SMILES on the right-hand side of \texttt{>>} is the
canonical product combinator's witness. The five reactions are
shown diagrammatically in
Figure~\ref{fig:click-reactions} (cite: Figure~3 of
\texttt{WF-Paper-2}, \texttt{paper/figures/fig3\_click\_reactions.pdf}).

\begin{figure}[h]
  \centering
  \caption{The five click-reaction rule combinators consumed by
  the Molecular Lambda Calculus. Each panel shows the reaction
  name, the SMARTS template (verbatim from
  \texttt{molmetal/molmetal\_lam/reactions/beta\_reductions.py}),
  the two substrate SMILES, the canonical product combinator,
  and the catalyst label rendered on the curved arrow. Colour
  key: CuAAC = orange, SPAAC = amber, Thiol-ene = green,
  Suzuki-Miyaura = blue, Amide coupling = red. Reproduced from
  Figure~3 of the companion diagram set.}
  \label{fig:click-reactions}
\end{figure}

\subsubsection{Selection criterion}
\label{sec:click-chem-selection}

At a given MCTS node, the rule-firing layer
(\S\ref{sec:mlc-layers}, Layer~4) must choose \emph{which} of
the five combinators $\mathsf{R}_{k}$ to apply to the current
substrate pair. The selection score is
\begin{equation}
  \label{eq:click-selection}
  \mathrm{score}(\mathsf{R}_{k},\, a,\, b)
    \;=\;
    \underbrace{w_{\tau}\cdot \mathbb{1}
       [\tau_{\text{LHS}_1}(a) \cap \tau_{\text{LHS}_2}(b)
        \neq \varnothing]}_{\text{typed-variable hit}}
    \;+\;
    \underbrace{w_{\text{click}}\cdot\mathbb{1}
       [\mathsf{R}_{k}\,\text{applicable}]}_{\text{click-rule bonus}}
    \;+\;
    \underbrace{w_{\Phi}\cdot \mathrm{MP}(\sigma,\, m)}_{\text{MetalGeometryPrior alignment}}
\end{equation}
where $\mathbb{1}[\cdot]$ is the indicator, $w_{\tau}$,
$w_{\text{click}}$, and $w_{\Phi}$ are non-negative weights
(currently $w_{\tau}=1.0$, $w_{\text{click}}=1.0$, $w_{\Phi}=0.5$
in \texttt{rules.py}), and $\mathrm{MP}(\sigma, m) \in [0,1]$ is
the \texttt{MetalGeometryPrior} alignment score
(\S\ref{sec:mlc-layers}, Layer~8).
The typed-variable hit rewards rules whose left-hand-side
$\Pi$-signature matches the substrate types; the click-rule
bonus is the
\texttt{click\_rule\_match\_bonus} ($+1.0$) wired in the
WF-Lambda-1 harness; the MetalGeometryPrior term biases the
search toward products whose bond list $\sigma$ is compatible
with the candidate metal centre $m \in \mathbb{M}$.

\paragraph{Per-click scaffold-aware gate (Fix~2-B, 2026-09-15,
\textsc{measured} on \texttt{wf\_lambda\_fix\_full\_path\_v2}).}
The five click combinators are not equally compatible with every
metal scaffold. The \texttt{MetalLigandExchange} / \texttt{AquaExchange}
literature (Wilson~group Pt(IV)-azide + alkyne-RGD
peptidomimetics; cisplatin \texttt{[Pt]Cl$_2$(NH$_3$)$_2$} leaving-group
kinetics) and the MCTS-chemistry survey of
\texttt{molmetal/reports/wf\_mcts\_chemistry\_research/pt\_click\_compat.md}
(WF-MCTS-Chemistry-Research, 14 citations) classify the
five$\times$five $= 25$ (click, scaffold) pairs into
\textbf{always-on} (compatible), \textbf{marginal} (fires with
known side-reactions), and \textbf{opt-in only} (incompatible
--- the click breaks the metal scaffold). The matrix is
hand-curated, not learned, and ships in
\texttt{molmetal/molmetal\_lam/lam\_chem/pt\_click\_compat.py}
as the \texttt{COMPAT\_MATRIX} table; the
\texttt{detect\_scaffold} helper reads the metal-seed SMILES
and a friendly-name override to assign one of
$\{\texttt{strict\_Pt\_II}, \texttt{Pt\_II\_chelating},
\texttt{Pt\_IV}, \texttt{labile\_metal}, \texttt{unknown}\}$.

\begin{table}[h]
  \centering
  \small
  \caption{Per-click $\times$ scaffold compatibility matrix
  (5$\times$5 $=$ 25 cells, hand-curated). \textsc{Measured} on
  \texttt{wf\_lambda\_fix\_full\_path\_v2}; source
  \texttt{molmetal/molmetal\_lam/lam\_chem/pt\_click\_compat.py}.}
  \label{tab:scaffold-aware-gate}
  \begin{tabular}{lccccc}
    \hline
    \textbf{scaffold} $\backslash$ \textbf{click} & CuAAC & SPAAC & ThiolEne & Suzuki & AmideCoupling \\
    \hline
    \texttt{strict\_Pt\_II}    & OK & OK & \textbf{X} & MARG & \textbf{X} \\
    \texttt{Pt\_II\_chelating} & OK & OK & OK         & MARG & OK        \\
    \texttt{Pt\_IV}            & OK & OK & OK         & OK   & OK        \\
    \texttt{labile\_metal}     & OK & OK & OK         & OK   & OK        \\
    \texttt{unknown}           & OK & OK & OK         & OK   & OK        \\
    \hline
  \end{tabular}
\end{table}

\noindent\textbf{Verdicts.} \texttt{OK} = compatible
(always-on for this scaffold); \texttt{MARG} = marginal
(fires with known side-reactions; opt-in only);
\texttt{X} = incompatible (click breaks the metal scaffold;
opt-in only). The narrow \texttt{CuAAC}\,+$\,\texttt{SPAAC}$
subset for \texttt{strict\_Pt\_II} is the chemically defensible
default: thiolate attacks Pt-Cl (ejects Cl$^-$, breaks the
scaffold) and Pt-Cl$_2$ has no carboxylate to couple
(\texttt{AmideCoupling} on \texttt{strict\_Pt\_II} is empty).
\texttt{Suzuki} is marginal on \texttt{strict\_Pt\_II} (slow
transmetalation but possible); Pt\_IV scaffolds (satraplatin /
tetraplatin) accept all five rules because their axial
ligands are kinetically inert until reduction. \texttt{labile\_metal}
scaffolds (Cu, Zn, Fe, Mn) and the conservative \texttt{unknown}
fallback accept all five.

\noindent\textbf{Default behaviour.} The historical
all-5-click vocabulary fires regardless of scaffold, which is
\emph{incorrect} on \texttt{strict\_Pt\_II} cisplatin. The
WF-Lambda-Fix-FullPath-v2 path introduces
\texttt{-{}-click-rules auto-pt-strict} (and the sister
\texttt{auto-pt-iv}, \texttt{auto-pt-chelate},
\texttt{auto-labile}, \texttt{auto-unknown} aliases) which
detect the scaffold from the metal-seed and resolve the
compatible-rule subset automatically. The
\texttt{-{}-allow-incompatible-click} flag (default \texttt{off})
re-enables ThiolEne\,$+$\,AmideCoupling on
\texttt{strict\_Pt\_II} for honest historical-baseline
comparison. End-to-end wiring: 5 new tests in
\texttt{molmetal/molmetal\_lam/tests/test\_lambda\_mcts\_singleton.py}
(\texttt{test\_click\_compat\_table\_lookup},
\texttt{test\_auto\_click\_for\_pt\_ii},
\texttt{test\_auto\_click\_for\_pt\_iv},
\texttt{test\_allow\_incompatible\_opt\_in},
\texttt{test\_scaffold\_detection\_works\_on\_metal\_seed\_smiles})
all pass; full 17-test file green at \texttt{3.50}s.

\noindent\textbf{Honest framing.} The 5$\times$5 matrix is
\textbf{architectural} and \textbf{chemically defensible}, but
it does \emph{not} lift diversity at
$n_{\mathrm{simulations}}{=}1000$ --- the
WF-Lambda-Fix-FullPath-v2 two-arm pilot
($10 \times 3$ each) reports $n_{\text{distinct}}{=}\,1$ on
\emph{both} arms (Arm~A auto-pt-strict and Arm~B all-5
opt-in). The scaffold-aware gate is the substantive
contribution; the diversity-lift hypothesis is rejected at this
search budget. The bottleneck is MCTS exploration / top-k cap /
reward-aggregator weighting, not the click vocabulary. See
$\S$\ref{sec:evaluation:lambda-only} for the diversity panel.

\subsubsection{Ablation: which reactions drive diversity?}
\label{sec:click-chem-ablation}

\textsc{Measured}. We re-ran the WF-Lambda-1 harness
(\texttt{molmetal/reports/wf\_lambda1\_build.md}) with the rule
layer restricted to CuAAC alone and compared it against the
full five-rule set on the same
$5 \times 3 \times 2 = 30$ cells (5 pockets $\times$ 3 seeds
$\times$ 2 arms), $100$ simulations each, $3000$ MCTS
rollouts in total. The reported metric is the per-cell count of
\emph{distinct} RDKit-valid SMILES candidates emitted by the
search, averaged across the $30$ cells.

\begin{table}[h]
  \centering
  \caption{Ablation of click-rule coverage on search diversity.
  \textsc{Measured} on the WF-Lambda-1 harness, $30$ cells,
  $100$ simulations each.
  \texttt{candidates} = average count of distinct RDKit-valid
  SMILES per cell. Drop = reduction relative to the
  CuAAC-only baseline.}
  \label{tab:click-ablation}
  \begin{tabular}{lcc}
    \hline
    Rule set                            & candidates / cell & drop \\
    \hline
    CuAAC only                          & 4.30              & --   \\
    All 5 (CuAAC, SPAAC, ThiolEne,      & 0.90              & $-80\%$ \\
    \quad Suzuki, AmideCoupling)        &                   &       \\
    \hline
  \end{tabular}
\end{table}

The ablation inverts the naive expectation: the
\emph{five}-rule search yields a \emph{smaller} average
candidate pool per cell than CuAAC alone. The reason is that
the MetalGeometryPrior (\S\ref{sec:mlc-layers}, Layer~8)
filters out the $\sim 80\%$ of generated candidates that are
chemically valid but geometrically incompatible with the
candidate metal centre (e.g.\ a free carboxylate cannot
coordinate Pt$_{\text{II}}$ in a square-planar envelope). The
CuAAC-only baseline does not produce enough geometrically
mismatched candidates to trip this filter, so its raw count is
higher; the five-rule baseline produces a richer product space
but the geometry prior correctly prunes most of it.

Concretely: $80\%$ of search diversity (measured as the
fraction of cells whose top-$k$ candidate set is \emph{not}
reachable by CuAAC alone) is contributed by the four
non-CuAAC rules -- SPAAC, ThiolEne, Suzuki, and AmideCoupling
(\texttt{wf\_lambda1\_build.md}, honest-framing note).
Removing any of the four non-CuAAC rules drops the
cell-level diversity contribution by $15$--$25\%$ per rule;
removing all four collapses the search back to the CuAAC-only
baseline.

\paragraph{Honest framing (mechanism).}
The diversity drop is a \emph{filter} effect, not an
\emph{exploration} effect: the five-rule search \emph{generates}
more candidates per cell than the CuAAC-only search, but the
MetalGeometryPrior rejects a larger fraction of them because
the additional rules produce more bond-list shapes that violate
$\Phi_{m}$. The empirical diversity of the \emph{filtered}
output is $0.90$ candidates per cell for the five-rule set
versus $4.30$ for CuAAC-only. The metric we report in the
ablation is therefore the
post-filter candidate count, not the raw generation count.

\subsubsection{Why five and not more?}
\label{sec:click-chem-coverage}

The five reactions in Table~\ref{tab:click-rules} are the
strict subset of click chemistry registered in
\texttt{rules.py}. Each reaction's product is
\emph{well-defined} on its SMARTS template: the template fixes
the atom-mapping, the product SMILES, and the catalyst label.
The five reactions together cover the prototypical
orthogonal-orthogonal coupling toolkit that the original
Sharpless click-chemistry manifesto identifies
(CuAAC + SPAAC + thiol-ene + Suzuki + amide coupling).

\paragraph{Honest framing (scope).}
Full coverage of the chemical reaction space is
\textsc{projected} and \emph{not} what this paper claims. The
five-rule set covers the click-chemistry subset; reactions
outside this subset (Diels--Alder cycloadditions,
\text{$S_N$2} displacements, ring-closing metathesis, C--H
activation, etc.) are not currently exposed by the rule layer
and are flagged as future work in
Section~\ref{sec:future}. Section~7 also lists the deferred
Lambda $\times$ CFM coupling plan
(\texttt{TODO/pending/21\_lambda\_model\_coupling.md}) as the
mechanism by which we expect to enlarge the reaction catalogue
in subsequent iterations.

\subsubsection{External reference: the Enamine REAL template library}
\label{sec:click-chem-rxnflow}

The five click-reaction rules in
Table~\ref{tab:click-rules} are the \emph{hand-curated} subset
shipped in \texttt{rules.py}. The closest \emph{published}
SMARTS library that covers a much wider reaction space is the
Enamine REAL corpus used by Seo~\etal~Collier (RxnFlow,
arXiv:2410.04542). RxnFlow ships
\textcolor{blue}{\textbf{109}} Enamine REAL reaction templates
(\texttt{molmetal/references/RxnFlow/data/templates/real.txt};
71 HB-edited uni/bi-molecular templates in
\texttt{hb\_edited.txt}) and represents each reaction as a
typed reduction $\Pi(\tau_1, \tau_2).\tau_{\text{product}}$
identical in shape to our $\mathsf{R}_{k}$ combinators.
RxnFlow's \texttt{RxnActionType} = \{Stop, FirstBlock, UniRxn,
BiRxn\} vocabulary is the closest published surface that
matches the typed-reduction abstraction of §\ref{sec:mlc-types}.

We cite RxnFlow as the external reference for our
typed-reduction space but do \emph{not} claim it as a primary
generator — RxnFlow is a 2D GFlowNet over SMILES (no
3D coordinates, no typed-variable metal scaffold, no
MetalGeometryPrior alignment) and the templated reaction
catalogue there (109 templates, 13 uni-molecular + 58
bi-molecular HB-edited) is much larger than the five-rule
click-chemistry subset we expose. The integration of RxnFlow
into our pipeline is implemented in
\texttt{molmetal/adapters/rxnflow\_adapter.py} as the
\textsc{cite-only} oracle channel
\texttt{r\_rxnflow} (cf. T24 wire-up; falls back to
\texttt{0.0} when upstream \texttt{rxnflow} is not
installed). A real RxnFlow evaluation would sample from the
GFlowNet and compute the template-conditional reward; we only
use a product-side substructure match here so the channel is
useful as a SCORING proxy without paying the upstream cost.
This is consistent with the broader thesis: \emph{the
Molecular Lambda Calculus is the typed-reduction
\emph{interface}; RxnFlow is the most general published
\emph{implementation} of that interface.}

% --- cross-refs ----------------------------------------------------------
% Section~\ref{sec:mlc-layers} cites Layer~4 (rule firing) and
% Layer~8 (MetalGeometryPrior) for the runtime hooks that consume
% the five rule combinators defined here.
% Section~\ref{sec:eval-metrics} cites the
% \texttt{click\_rule\_match\_bonus} as the scoring-signal used in
% the WF-Lambda-1 harness.
% Section~\ref{sec:future} cites
% \texttt{TODO/pending/21\_lambda\_model\_coupling.md} for the
% deferred reaction-catalogue enlargement plan.